% ==================================================
% MACROS - LENH TU DINH NGHIA
% ==================================================

% ==================================================
% BIEN QUAN LY DE THI
% ==================================================
\newcommand{\tendonvi}{DHPX PREP}
\newcommand{\namthi}{2027}
\newcommand{\monthi}{TOÁN}
\newcommand{\thoigianlambai}{90 phút}
\newcommand{\sotrangde}{0x}

\newcommand{\sode}{1}
\newcommand{\tendedethi}{ĐỀ ÔN LUYỆN HÀM SỐ}

% =========================
% MA DE
% =========================
\newcommand{\made}{0101}

% =========================
% CHE DO GIAO VIEN / HOC SINH
% =========================
\newif\ifgiaovien

% Bat loi giai:
\giaovientrue

% An loi giai:
%\giaovienfalse

% Moi truong loi giai an/hien an toan
\NewEnviron{loigiai}{%
  \ifgiaovien
    \par\noindent\textbf{Lời giải.}\par\noindent
    \BODY
    \par
  \fi
}

% =========================
% CAU HOI
% =========================
\newcounter{cau}

\newcommand{\cau}{%
  \stepcounter{cau}%
  \par\vspace{3pt}%
  \noindent\textbf{Câu \thecau. }%
}

% =========================
% CAU TRA LOI TRAC NGHIEM
% =========================
\settasks{
  label=\textbf{\Alph*.},
  label-width=18pt,
  item-indent=22pt,
  label-offset=4pt,
  column-sep=10pt,
  after-item-skip=4pt,
  after-skip=6pt
}

\newcommand{\bonpa}[4]{%
\begin{tasks}(4)
  \task #1
  \task #2
  \task #3
  \task #4
\end{tasks}
}

\newcommand{\haipa}[4]{%
\begin{tasks}(2)
  \task #1
  \task #2
  \task #3
  \task #4
\end{tasks}
}

\newcommand{\motcot}[4]{%
\begin{tasks}(1)
  \task #1
  \task #2
  \task #3
  \task #4
\end{tasks}
}

% =========================
% PHAN DE THI
% =========================
\newcommand{\tenphanhientai}{}
\newcommand{\phan}[2]{%
  \vspace{6pt}
  \renewcommand{\tenphanhientai}{#1}%
  \par\noindent\textbf{PHẦN #1.} #2
  \vspace{3pt}
}

% =========================
% DUNG / SAI
% =========================
\newenvironment{dungsai}
{%
  \begin{enumerate}[
    label=\textbf{\alph*)},
    leftmargin=18pt,
    itemsep=2pt,
    topsep=2pt
  ]
}
{%
  \end{enumerate}
}

% =========================
% TRA LOI NGAN
% =========================
%\newcommand{\shortanswerbox}[1]{%
  %\hfill
  %\fbox{%
    %\parbox[c][0.6cm][c]{2.8cm}{%
      %\centering
     % \ifgiaovien
        %#1%
    %  \else
    %    \phantom{#1}%
    %  \fi
%    }%
%  }%
%}

\newcommand{\shortanswerline}[1]{%
  \hfill
  \ifgiaovien
    \underline{\hspace{0.3cm}#1\hspace{0.3cm}}%
  \else
    \underline{\hspace{3cm}}%
  \fi
}

\newcommand{\fourshortanswers}[4]{%
\begin{center}
\begin{tabular}{cccc}
\textbf{a)} & \textbf{b)} & \textbf{c)} & \textbf{d)} \\[4pt]
\fbox{\parbox[c][0.6cm][c]{2.2cm}{\centering \ifgiaovien #1\else\phantom{#1}\fi}} &
\fbox{\parbox[c][0.6cm][c]{2.2cm}{\centering \ifgiaovien #2\else\phantom{#2}\fi}} &
\fbox{\parbox[c][0.6cm][c]{2.2cm}{\centering \ifgiaovien #3\else\phantom{#3}\fi}} &
\fbox{\parbox[c][0.6cm][c]{2.2cm}{\centering \ifgiaovien #4\else\phantom{#4}\fi}}
\end{tabular}
\end{center}
}

% =========================
% O MA DE
% =========================
\newcommand{\omade}{%
\begin{tikzpicture}
\node[
  draw,
  minimum width=31mm,
  minimum height=8mm,
  inner sep=0pt
]
{\textbf{Mã đề: \made}};
\end{tikzpicture}
}

% =========================
% HEADER DE THI
% =========================
\newcommand{\tieudethi}{%
\begin{center}
\begin{tabularx}{\textwidth}{>{\centering\arraybackslash}X>{\centering\arraybackslash}X}
\textbf{\tendonvi}
&
\textbf{KỲ THI TỐT NGHIỆP TRUNG HỌC PHỔ THÔNG NĂM \namthi}
\\
\textbf{\tendedethi}
&
\textbf{Môn thi: \monthi}
\\
\textit{(Đề thi có \sotrangde\ trang)}
&
\textit{Thời gian làm bài \thoigianlambai, không kể thời gian phát đề}
\end{tabularx}
\end{center}

\vspace{4pt}

\begin{tabularx}{\textwidth}{X r}
\textbf{Họ, tên thí sinh:} \dotfill & \multirow{2}{*}{\omade} \\
\textbf{Số báo danh:} \dotfill &
\end{tabularx}

\vspace{8pt}
}

% =========================
% LENH TIEN ICH
% =========================
\newcommand{\chon}[1]{\textbf{Chọn #1.}}
\newcommand{\dapanso}[1]{Đáp án: #1.}
\newcommand{\hoac}{\quad\text{hoặc}\quad}

\newcommand{\R}{\mathbb{R}}
\newcommand{\N}{\mathbb{N}}
\newcommand{\Z}{\mathbb{Z}}
\newcommand{\Q}{\mathbb{Q}}
\newcommand{\C}{\mathbb{C}}

\newcommand{\vect}[1]{\overrightarrow{#1}}
\newcommand{\dd}{\,\mathrm{d}}

\newcommand{\abs}[1]{\left|#1\right|}
\newcommand{\paren}[1]{\left(#1\right)}
\newcommand{\bracket}[1]{\left[#1\right]}
\newcommand{\set}[1]{\left\{#1\right\}}







% =========== % Lời giải ví dụ % ================ % 


% Lời giải ví dụ này sẽ có cái ô màu xanh khá là chảnh chó ====))))) % 


% ==================================================
% LOI GIAI TOI GIAN - ZEN STYLE
% ==================================================

% =========================
% LOI GIAI CHO VI DU MINH HOA
% Luon hien, vi day la phan ly thuyet - vi du
% =========================
\NewEnviron{loigiaividu}{%
  \par\noindent\textbf{Lời giải ví dụ.}\par\noindent
  \BODY
  \par
}

% =========================
% LOI GIAI THAM KHAO CHO BAI TAP TONG HOP
% An/hien theo che do giao vien
% =========================
\NewEnviron{loigiaithamkhao}{%
  \ifgiaovien
    \par\noindent\textbf{Lời giải.}\par\noindent
    \BODY
    \par
  \fi
}

% =========================
% DANH SACH LUU LOI GIAI THAM KHAO
% =========================
\newcommand{\danhsachloigiaithamkhao}{}
\newcommand{\danhsachdapan}{}

\newcounter{cauLG}

\newcommand{\xoaloigiaithamkhao}{%
  \global\let\danhsachloigiaithamkhao\empty
  \global\let\danhsachdapan\empty
  \setcounter{cauLG}{0}
}

% =========================
% CAU HOI CO LOI GIAI GOM VE CUOI
% #1: De bai
% #2: Loi giai
% =========================
\newcommand{\caugiai}[2]{%
  \cau #1%
  \ifgiaovien
    \gappto\danhsachloigiaithamkhao{%
      \stepcounter{cauLG}%
      \par\vspace{6pt}%
      \noindent\textbf{Câu \thecauLG. }#1\par
      \begin{loigiaithamkhao}
      #2
      \end{loigiaithamkhao}
    }%
  \fi
}

% =========================
% IN TOAN BO LOI GIAI THAM KHAO
% =========================
\newcommand{\inloigiaithamkhao}{%
  \ifgiaovien
    \setcounter{cauLG}{0}%
    \danhsachloigiaithamkhao
  \fi
}

% =========================
% QUAN LY DE THI (BAT DAU / KET THUC)
% =========================
\newcommand{\batdaude}[3]{%
  % #1: So de, #2: Ten de, #3: Ma de
  \renewcommand{\sode}{#1}%
  \renewcommand{\tendedethi}{#2}%
  \renewcommand{\made}{#3}%
  \newpage
  \setcounter{cau}{0}%
  \xoaloigiaithamkhao
  \tieudethi
}

\newcommand{\ketthucde}{%
  \ifgiaovien
    \newpage
    \begin{center}
      \LARGE\textbf{ĐÁP ÁN VÀ LỜI GIẢI CHI TIẾT ĐỀ \sode}
    \end{center}
    \inloigiaithamkhao
  \fi
}